\clearpage
\subsection{Compatibility Contract}

This section defines the default compatibility contract for reserved fields and optional features. A later section
may narrow one of these defaults for a specific structure, but it must say so explicitly.

This manual treats reserved and unallocated encodings as one architectural category: software must not depend on their
value, decode result, or behavior.

\subsection{Reserved Fields}
Architected registers and architectural memory structures use these defaults unless their defining section gives a
narrower rule.

Software-defined fields are not reserved fields. Software may store values in them, and hardware ignores them except
where the containing structure explicitly says otherwise.

\Needspace{1.25in}
\manualtablecaption{Reserved Field Defaults}
\begingroup\footnotesize
\setlength{\aboverulesep}{0pt}
\setlength{\belowrulesep}{0pt}
\setlength{\extrarowheight}{1.2pt}
\begin{center}
\begin{tabularx}{\linewidth}{@{}>{\raggedright\arraybackslash}p{1.65in}>{\raggedright\arraybackslash}p{0.75in}>{\raggedright\arraybackslash}p{1.35in}>{\raggedright\arraybackslash}X@{}}
\toprule
\rowcolor{ManualHeaderFill}
\textbf{Field Class} & \textbf{Read Rule} & \textbf{Write or Use Rule} & \textbf{Fault}\\
\midrule
Architected register reserved bits & zero & must be zero & --\\
Control-register reserved bits & zero & must be zero & \texttt{INVALID\_CONTROL\_STATE}\\
Reserved selector values & -- & invalid selector & \texttt{INVALID\_CONTROL\_STATE}\\
Consumed reserved PTE bits & -- & invalid translation input & \texttt{PAGE\_FAULT}\\
Reserved event-code classes or IDs & -- & must not be generated by the base profile & --\\
Reserved architectural event-frame bits & -- & must be zero & \texttt{INVALID\_CONTROL\_STATE} on ERET\\
\bottomrule
\end{tabularx}
\end{center}
\endgroup


\subsection{Reserved Encodings}
Reserved instruction opcodes, reserved extension opcodes, reserved effective-address forms, and optional
instruction-group encodings whose CPUID feature bit is clear raise \texttt{ILLEGAL\_INSTRUCTION} during decode.

Noncanonical encodings are invalid unless an instruction description explicitly defines an alias or priority rule.
Assemblers emit canonical encodings by default; disassemblers prefer canonical spellings.

\Needspace{3.6in}
\subsection{CPUID Compatibility}
CPUID is the compatibility boundary for optional architectural behavior. Software must not use an optional
instruction group, optional register state, or optional state image unless the corresponding CPUID bit or leaf
reports support.

Unknown selectors return zero. Software ignores reserved CPUID result bits. CPUID is unprivileged. For a logical
processor, CPUID results remain stable until that logical processor is reset.
